%%%%%%%%%%%%%%%%%%%%%%%%%%%%%%%%%%%%%%%%%%%%%%%%%%%%%%%%%%%%%%%%%%%%%%%%%%%%%%%%
%%
\cleardoubleoddpage%  Make sure to start each chapter on a new odd page

\chapter{Solver Relation and Framework}
\label{chapter:solver-relation-framework}
\ac{SAT}, \ac{SMT} and \ac{MaxSAT} (as extensions of \ac{SAT}) and \ac{ILP} represent four major paradigms for constraint solving in computer science and optimization. Their differences in expressiveness, solver technology, and application domains justify their inclusion in a comprehensive study of automated reasoning approaches. Especially with \ac{SAT} being one of the most popular paradigms for constraint solving~\cite{bib:2023-ulrich-oltean}, while \ac{SMT} is a required extension to properly cover our language, and \ac{ILP} being a good fit for our problem domain, we will focus on these four paradigms in the following sections.

\section{Logical Solvers}
The four following chapters will introduce the most basic concept of the respective solver paradigms. The field of possible constraint solvers is vast and far from complete. We restrict ourselves to the four most prominent options, to ensure the focus of the thesis stays on the core concept of our \Ac{IR}.

\subsection{Boolean Satisfiability Problem (SAT)}
It is proven that \Ac{SAT} is NP-complete~\cite{bib:2008-marques}. It means that if a problem can be reduced to \ac{SAT} it can be solved by a \ac{SAT} solver with at worst-case exponential runtime complexity. Using this approach, we handle propositional logic formulas, which are representations of boolean variables possibly combined with logical operators like \textbf{AND} ($\land$), \textbf{OR} ($\lor$), and \textbf{NOT} ($\lnot$). Most solvers are specialized on formulas in \Ac{CNF}, which means there is a conjunction ($\land$) of clauses, which in turn is a disjunction ($\lor$) of literals (boolean variable or their negation). An example formula in \ac{CNF} would be: $a \land (b \lor \lnot c)$. While \ac{SAT} is also powerful enough to express linear relations in a \ac{CNF} it is not the most efficient way to do so~\cite{bib:2008-marques}. This is why we take a look at \ac{SMT} in the next section.

\subsection{Satisfiability Modulo Theories (SMT)}
\Ac{SMT} is one of the most promising extensions of \ac{SAT}~\cite{bib:2008-marques} hence, we also handle it in this chapter. In particular \ac{SMT} is the middle ground between first or higher-order-logic and \ac{SAT}~\cite{bib:2018-clarke}. It tries to mend the gap between expressiveness and efficiency of automatically solving problems. By introducing the likes of equality, arithmetics and uninterpreted functions we get a more powerful language, which makes expressing problems easier~\cite{bib:2007-demoura}. On the other hand, this expressiveness comes with a trade-off in solver efficiency and hence, we need to introduce specialized theory solvers. Similarly to \ac{SAT}, \ac{SMT} is the conjunction, or disjunction of formulas (or again their negations). An example formula in \ac{SMT} would be: $\Phi = x + y \leq 10 \land (z = x * 2 \lor \lnot (y < 5)) \land f(x) = 4$. Additionally, we define $M$ to be an assignment of the formula $\Phi$, and state that $M \models \Phi$ if $M$ satisfies $\Phi$. In \ac{SMT} we see the usage of arithmetics, which is not possible in \ac{SAT}. This makes \ac{SMT} a more expressive language, which is easier to use for humans. Similar to \ac{SAT}, \ac{SMT} can also express linear relations, but again it is not the most efficient way to do so~\cite{bib:2008-marques}. As \ac{SMT} comes with a cost of efficiency, we will also take a look at \ac{MaxSAT} in the next section, which does not have this drawback with the downside of being as expressive as \ac{SAT}.

\subsection{Maximum Satisfiability Problem (MaxSAT)}
We already covered \ac{SAT} and extension of it in form of \ac{SMT}. Another prominent extension of \ac{SAT} is \Ac{MaxSAT}~\cite{bib:2013-ansotegui}. Similarly to \ac{SAT} it consists of boolean variables combined using disjunction $\left( \bigvee_{i=1}^N p_i \right)$ --- so-called clauses. In turn these clauses are combined using conjunction $\left( \bigwedge_{j=1}^M c_j \right)$. In contrast to basic \ac{SAT}, \ac{MaxSAT} has the addition of soft clauses. This is a good fit for our definition of the problem, as we also have hard and soft constraints. We define a set of weighted clauses~\cite{bib:2013-ansotegui}
$$\phi = \{(C_1, w_1), ... (C_m, w_m)\}$$
If one wants to define hard constraint-clause, one has to simply set $w = \infty$. The goal of \ac{MaxSAT} is to find an assignment for $\phi$ with as little cost as possible. This is also a perfect fit for our problem definition, as we want an assignment with as little cost as possible, while still satisfying all hard constraints. \\
As already mentioned for \ac{SMT} our language mostly consists of linear relations and inequalities. While \ac{MaxSAT} can express these relations, it is not the most straight forward way~\cite{bib:2006-nieuwenhuis}. Therefore, we note that there is also a combination of \ac{MaxSAT} and \ac{SMT} called \Ac{MaxSMT}, which is basically weighted \ac{SMT}~\cite{bib:2016-pasareanu}. Also, as arithmetics is what is needed to express the majority of constraints we defined in our language. As such, we will also take a look at \ac{ILP} in the next section, as it is specialized on solving linear relations.

\subsection{Integer Linear Programming (ILP)}
\Ac{ILP} is designed to either minimize or maximize an objective function, while holding inequality and equality constraints or keeping integrality restrictions on variables~\cite{bib:1994-chaudhuri}. Similarly to \ac{SAT} we know that \ac{ILP} is NP-complete. This means that if a problem can be reduced to \ac{ILP} it can be solved by an \ac{ILP} solver with at worst-case exponential runtime complexity. An example of an \ac{ILP} problem would be: $z = \min{(c^Tx | Ax \leq b, x \in \Z^n)}$. Here, $c^Tx$ is the objective function to minimize and $z$ represents the corresponding minimum value. $Ax \leq b$ are the constraints, and $x \in \Z^n$ enforces integrality on the variables. As we can see, \ac{ILP} is specialized on linear relations, which makes it a good fit for our problem domain. Additionally, we have the issue of minimizing cost, which is another point in favor of \ac{ILP} solvers but is also handled in \ac{MaxSAT}.

\subsubsection{Example Encoding}
As most of our constraint focus on linear inequality relations ($a + b \leq c$) and minimizing cost ($\min{(c(a))}$) we decided to use \Ac{ILP} as our showcase example. Hence, suggest the following constraints as examples:
\begin{table}[!ht]
    \centering
    \begin{tabular}{|c|c|}
        \hline
        Constraint                          & \ac{ILP}-like Encoding                                                                                                                           \\
        \hline
        Task A before Task B                & $\operatorname{a}(A) + \operatorname{dur}(A) \leq \operatorname{a}(B)$                                                                           \\
        \hline
        Task A with interval [s, e]         & $s \leq \operatorname{a}(A) \land \operatorname{a}(A) + \operatorname{dur}(A) \leq e$                                                            \\
        \hline
        Task A must not overlap with Task B & $\operatorname{a}(A) + \operatorname{dur}(A) \leq \operatorname{a}(B) \lor \operatorname{a}(B) + \operatorname{dur}(B) \leq \operatorname{a}(A)$ \\
        \hline
        Minimize cost                       & $\min{(c(a))}$                                                                                                                                   \\
        \hline
    \end{tabular}
    \caption{Constraint conversion table to \ac{ILP}-like constraints.}
\end{table}\\
In this case we inserted the function $\operatorname{a}$ instead of using variables to showcase that the output of such a solver would be an assignment of time slots (start times) to each task. It becomes clear that the constraints we defined in our language can be converted quite naturally to \ac{ILP} constraints or at least a form of inequalities. As for the last constraint of minimizing cost, we can instead of looking at the minimization of the whole assignment, also split this with the definition of the cost function on a specific task. Therefore, we can see that we do not need to define a complete assignment which is minimized, but we minimize the cost of each task itself.

\subsection{Solver comparison}
In the previous sections we introduced for major solver paradigms for constraint solving. As minimizing cost while holding both hard and soft constraints is a major part of the problem definition, we can see that both \ac{MaxSAT} and \ac{ILP} are good fits for our problem at hand. However, determining which one is the better fit is not a straight forward task~\cite{bib:2013-ansotegui}. It is even not clear if there is one best solver between those. It can even be beneficial to combine both paradigms to get the best performance~\cite{bib:2025-zhang}. Meaning, An \ac{ILP} solver could be used as a preprocessing step to reduce the problem to be later given to a \ac{MaxSAT} solver. Therefore, we conclude that from the four mentioned solver types, we can not determine one best solver. Instead, we designed our \Ac{IR} to be as solver-agnostic as possible, to allow for future research to determine the best fitting solver or combination of solvers for our specific use-case.

%% 
%%%%%%%%%%%%%%%%%%%%%%%%%%%%%%%%%%%%%%%%%%%%%%%%%%%%%%%%%%%%%%%%%%%%%%%%%%%%%%%%